\section{为什么需要工作底稿}
光通信产业链横跨材料、半导体设备、光电芯片、精密器件、模块、光纤光缆、网络设备和下游应用。若只写市场叙事，很容易把 AI 模块的高景气错误外推到所有层级；若只写公司财务，又会忽视上游认证、良率、标准演进和应用周期。本附录把后续调研拆成可执行工作底稿，用于每次更新报告时逐项核查。

\section{全链条调研问题库}
\begin{exhibitbox}[表：产业链调研问题库]
\centering\tiny
\renewcommand{\arraystretch}{1.20}
\setlength{\tabcolsep}{2pt}
\begin{tabularx}{\exhibitboxwidth}{>{\bfseries\raggedright\arraybackslash}p{1.35cm} >{\raggedright\arraybackslash\sloppy\hspace{0pt}}p{2.55cm} >{\raggedright\arraybackslash\sloppy\hspace{0pt}}p{2.45cm} >{\raggedright\arraybackslash\sloppy\hspace{0pt}}p{2.35cm} >{\raggedright\arraybackslash\sloppy\hspace{0pt}}X}
\toprule
\textbf{层级} & \textbf{核心问题} & \textbf{量化指标} & \textbf{优先来源} & \textbf{估值影响} \\
\midrule
材料/基底 & 预制棒、InP/GaAs、硅光晶圆、薄膜铌酸锂是否进入稳定量产 & 良率、缺陷密度、客户认证、扩产资本开支 & 公司公告、招股书、产业标准、客户验证记录 & 影响上游芯片/器件长期 PE 上限，不直接推高模块 EPS \\
制造设备 & 外延、主动耦合、老化和高速测试能否支撑 1.6T/CPO 放量 & 新签订单、交付周期、毛利率、费用率、现金流 & 公司订单公告、客户扩产、设备招标、调研纪要 & 亏损或项目制公司不纳入 PE 目标价，进入观察池 \\
光芯片/光源 & EML/CW/VCSEL/PD/APD/PLC/AWG 是否通过头部客户认证 & 客户数、收入占比、毛利率、研发费用率、良率 & 财报、客户公告、供应链验证、专利与产品手册 & 通过认证后提升倍数；若停在送样阶段则下修 \\
光器件/无源 & 高速模块 attach rate 是否提升，是否被客户自供替代 & 单季收入、毛利率、客户集中、存货、OCF & 季报、调研、模块厂需求、产品拆解 & attach rate 上升支持溢价；价格压力降低目标 PE \\
光模块 & 800G 是否持续放量，1.6T 是否从认证转批量 & 收入增速、毛利率、ASP、库存、应收账款、OCF & 季报、客户 capex、NVIDIA/云厂商动向、行业预测 & 直接决定 EPS 和目标价，是全链条 beta 核心 \\
光纤光缆/ODN & 运营商、海缆、FTTH、DCI 项目是否改善 & 招标量、项目毛利率、预制棒利用率、光纤价格、现金流 & 运营商招标、公司公告、行业价格、项目交付 & 慢周期底盘，按现金流和项目节奏给折价倍数 \\
网络设备 & OTN/ROADM/PON/交换路由产品是否改善利润率 & 运营商 capex、订单、毛利率、海外收入、研发投入 & 运营商招标、设备商财报、标准组织、云网络产品发布 & 更稳但纯度低，除非 AI 网络产品兑现否则不套模块倍数 \\
下游应用 & AI、DCI、运营商、FTTH、企业/工业是否同步兑现 & 云 capex、端口速率、招标、项目交付、行业库存 & 云厂商/芯片厂公告、运营商招标、行业协会、财报 & 决定估值节奏：AI 快，运营商慢，工业/汽车作为长期期权 \\
\bottomrule
\end{tabularx}
\end{exhibitbox}

\section{更新频率与触发规则}
\begin{exhibitbox}[表：更新触发规则]
\centering\scriptsize
\begin{tabularx}{\exhibitboxwidth}{>{\bfseries\raggedright\arraybackslash}p{2.0cm} >{\raggedright\arraybackslash\sloppy\hspace{0pt}}X >{\raggedright\arraybackslash\sloppy\hspace{0pt}}X}
\toprule
\textbf{频率} & \textbf{必须更新的内容} & \textbf{触发原因} \\
\midrule
每日/行情异常 & 当前价、市值、目标价空间、成交额、涨跌幅和拥挤度 & 光通信板块波动大，旧目标价空间很快失真 \\
月度 & 云厂商 capex、交换 ASIC/800G/1.6T 新闻、运营商招标、光纤价格 & 订单和标准变化通常领先财报 \\
季度 & 25 只覆盖标的收入、利润、EPS、毛利率、OCF、存货、应收账款 & 估值分母必须随财报重算 \\
事件触发 & 1.6T 批量订单、CPO/硅光量产、重大客户认证、出口管制、并购扩产 & 改变估值倍数或失效条件 \\
报告复审 & 来源注册表、claim audit、估值审计、PDF 可读性和 verifier & 防止数据过时、章节空心化和估值缺失 \\
\bottomrule
\end{tabularx}
\end{exhibitbox}

\section{上修与下修的硬门槛}
上修不能只因为行业新闻热。至少需要满足以下三类证据之一：第一，财报层面出现收入、毛利率和经营现金流同步改善；第二，客户认证从送样进入批量订单，并且公司能披露或通过利润表验证；第三，应用层需求从单一 AI 模块外溢到 DCI、运营商、光纤光缆或网络设备项目。若只有股价上涨而没有上述证据，本报告只更新当前价和空间，不上修目标 PE。

下修同样需要规则化。若 Q2/Q3 净利润低于 Q1 的 0.85 倍、经营现金流连续弱于利润、毛利率明显下滑、存货或应收账款快于收入增长、光纤光缆价格继续下行、运营商 capex 延后、或 CPO 改变可插拔模块价值分配，必须下调 EPS、PE 或动作标签。对于观察池公司，若亏损扩大或订单兑现继续延后，维持观察池并不得进入最终估值表。

\section{分层数据源优先级}
\begin{exhibitbox}[表：数据源优先级]
\centering\scriptsize
\begin{tabularx}{\exhibitboxwidth}{>{\bfseries\raggedright\arraybackslash}p{1.5cm} >{\raggedright\arraybackslash\sloppy\hspace{0pt}}p{3.0cm} >{\raggedright\arraybackslash\sloppy\hspace{0pt}}p{3.0cm} >{\raggedright\arraybackslash\sloppy\hspace{0pt}}X}
\toprule
\textbf{优先级} & \textbf{可用来源} & \textbf{用途} & \textbf{限制} \\
\midrule
A & 公司公告、财报、监管披露、结构化财务包 & 收入、利润、EPS、现金流、股本、目标价分母 & 滞后于订单和技术变化 \\
A- & 官方公司新闻、NVIDIA/Broadcom/OIF/Ethernet Alliance 等官方公开页 & 技术路线、标准、需求方向和供应链关系 & 通常不披露 A 股公司直接订单金额 \\
B & 行业机构公开摘要、公开演讲、招标公告、价格跟踪 & 需求节奏、应用周期、行业景气度 & 需要避免把摘要当作完整预测模型 \\
C & 二手媒体、社交平台、未取得原文的卖方片段 & 识别市场情绪和分歧 & 不得作为估值输入，不得替代目标价模型 \\
\bottomrule
\end{tabularx}
\end{exhibitbox}

\section{后续调研排期}
下一轮更新优先级如下。第一，逐只补齐 25 家覆盖公司的正式 2026Q1/半年度公告原文归档，减少第三方公告入口依赖。第二，单独建立光芯片/光源专题，核查源杰科技、长光华芯、仕佳光子在客户认证和良率上的真实进度。第三，补光纤光缆价格和运营商招标数据库，用于校准长飞、亨通、中天、通鼎、永鼎的慢周期估值。第四，跟踪中兴、烽火、光迅、锐捷等在 OTN、ROADM、PON、云网络设备上的利润率变化。第五，若罗博特科或其他设备链标的恢复盈利，再建立非 PE 或 EV/Sales/订单覆盖倍数模型。


\section{逐公司调研底稿}
本节不是投资建议重复，而是后续复盘时的逐公司核查模板。每家公司都需要同时检查产业位置、财务分母、估值结论、催化剂、失效条件和下一轮数据更新点。若后续行情或财报变化，只要替换当前价和 EPS 分母，就能快速判断动作标签是否需要改变。

\subsection{中际旭创（300308）}
中际旭创 的产业定位是 \textbf{AI 数据中心高速光模块}，在本报告中的层级为 核心光模块，研究权重 13\%。当前价 CNY1253.89，按 2026Q1 归母净利和 EPS 反推的市值约 CNY13881 亿元。2026Q1 收入 CNY194.96 亿元，归母净利 CNY57.35 亿元，毛利率 46.1\%，经营现金流 CNY33.68 亿元。以上是估值分母，不允许被行业叙事替代。

估值上，本报告采用 \textbf{PE/PEG} 方法：基于正常化 EPS 的 forward PE/PEG，用于盈利兑现强且订单持续的 AI 光模块龙头。中际旭创 的 2026E/2027E/2028E EPS 分别为 CNY22.52、CNY29.73、CNY36.27，2027E 收入/股为 CNY101.07，BVPS 为 CNY32.03；内在组件权重为 PE 100\%，内在价值锚 CNY1337.79，市场情绪锚 CNY1003.11，综合权重 内在65\% / 市场20\% / 券商15\%，综合目标价 CNY1229.56，综合空间 -1.9\%，动作 中性，证据质量 A-。二级校验为：客户/订单持续性、毛利率保持、PEG 排名和 PS 交叉校验。

下一轮调研首先验证催化剂：1.6T 放量、海外云厂商订单能见度提升、毛利率维持在 42\%以上。 其次验证失效条件：800G/1.6T 订单放缓、客户集中度压力上升，或毛利率跌破 38\%。 若催化剂没有进入收入、毛利率和现金流，不能上调 PE；若失效条件出现，应同时下调 EPS 和动作标签。对于产业链映射，必须核查它的增长是否来自本层级真实需求，而不是来自板块情绪扩散。


\subsection{新易盛（300502）}
新易盛 的产业定位是 \textbf{高速光模块}，在本报告中的层级为 核心光模块，研究权重 12\%。当前价 CNY566.00，按 2026Q1 归母净利和 EPS 反推的市值约 CNY5620 亿元。2026Q1 收入 CNY83.38 亿元，归母净利 CNY27.80 亿元，毛利率 49.2\%，经营现金流 CNY6.84 亿元。以上是估值分母，不允许被行业叙事替代。

估值上，本报告采用 \textbf{PE/PEG} 方法：基于正常化 EPS 的 forward PE/PEG，用于盈利兑现强且订单持续的 AI 光模块龙头。新易盛 的 2026E/2027E/2028E EPS 分别为 CNY12.17、CNY15.83、CNY19.31，2027E 收入/股为 CNY47.46，BVPS 为 CNY20.50；内在组件权重为 PE 100\%，内在价值锚 CNY633.04，市场情绪锚 CNY452.80，综合权重 内在65\% / 市场20\% / 券商15\%，综合目标价 CNY576.55，综合空间 +1.9\%，动作 中性，证据质量 A-。二级校验为：客户/订单持续性、毛利率保持、PEG 排名和 PS 交叉校验。

下一轮调研首先验证催化剂：海外客户订单延续、1.6T 产品认证、现金转化改善。 其次验证失效条件：订单前置后回落、模块 ASP 下跌超过 20\%，或经营现金流明显落后利润。 若催化剂没有进入收入、毛利率和现金流，不能上调 PE；若失效条件出现，应同时下调 EPS 和动作标签。对于产业链映射，必须核查它的增长是否来自本层级真实需求，而不是来自板块情绪扩散。


\subsection{天孚通信（300394）}
天孚通信 的产业定位是 \textbf{精密光器件与封装平台}，在本报告中的层级为 精密光器件，研究权重 6\%。当前价 CNY318.00，按 2026Q1 归母净利和 EPS 反推的市值约 CNY2472 亿元。2026Q1 收入 CNY13.30 亿元，归母净利 CNY4.92 亿元，毛利率 56.6\%，经营现金流 CNY1.82 亿元。以上是估值分母，不允许被行业叙事替代。

估值上，本报告采用 \textbf{PE/PEG} 方法：基于正常化 EPS 的 forward PE/PEG，用于盈利兑现强且订单持续的 AI 光模块龙头。天孚通信 的 2026E/2027E/2028E EPS 分别为 CNY2.64、CNY3.27、CNY3.86，2027E 收入/股为 CNY8.84，BVPS 为 CNY7.69；内在组件权重为 PE 100\%，内在价值锚 CNY147.20，市场情绪锚 CNY254.40，综合权重 内在57\% / 市场28\% / 券商15\%，综合目标价 CNY254.40，综合空间 -20.0\%，动作 中性观察，证据质量 B+。二级校验为：客户/订单持续性、毛利率保持、PEG 排名和 PS 交叉校验。

下一轮调研首先验证催化剂：微光学器件和封装件在 1.6T/CPO 产品中的附加值提升。 其次验证失效条件：器件降价、客户自供比例提高，或毛利率回落至 50\%以下。 若催化剂没有进入收入、毛利率和现金流，不能上调 PE；若失效条件出现，应同时下调 EPS 和动作标签。对于产业链映射，必须核查它的增长是否来自本层级真实需求，而不是来自板块情绪扩散。


\subsection{光迅科技（002281）}
光迅科技 的产业定位是 \textbf{电信/数通光器件与光模块}，在本报告中的层级为 光器件/光模块，研究权重 4\%。当前价 CNY240.34，按 2026Q1 归母净利和 EPS 反推的市值约 CNY1922 亿元。2026Q1 收入 CNY27.73 亿元，归母净利 CNY2.40 亿元，毛利率 26.8\%，经营现金流 CNY1.96 亿元。以上是估值分母，不允许被行业叙事替代。

估值上，本报告采用 \textbf{PE/PB/PS} 方法：70\% PE + 10\% PB + 20\% PS，用于有正 EPS 但 AI 纯度不完美的模块/器件平台。光迅科技 的 2026E/2027E/2028E EPS 分别为 CNY1.36、CNY1.64、CNY1.88，2027E 收入/股为 CNY18.91，BVPS 为 CNY12.82；内在组件权重为 PE 70\% / PB 10\% / PS 20\%，内在价值锚 CNY58.46，市场情绪锚 CNY187.47，综合权重 内在47\% / 市场38\% / 券商15\%，综合目标价 CNY187.47，综合空间 -22.0\%，动作 中性观察，证据质量 B。二级校验为：数通收入占比、毛利率桥、收入规模和现金转化。

下一轮调研首先验证催化剂：数通收入占比提高、硅光/CPO 产品验证推进。 其次验证失效条件：电信周期走弱、数通结构未改善，或毛利率持续低于 28\%。 若催化剂没有进入收入、毛利率和现金流，不能上调 PE；若失效条件出现，应同时下调 EPS 和动作标签。对于产业链映射，必须核查它的增长是否来自本层级真实需求，而不是来自板块情绪扩散。


\subsection{华工科技（000988）}
华工科技 的产业定位是 \textbf{光模块与激光设备混合平台}，在本报告中的层级为 混合光通信平台，研究权重 4\%。当前价 CNY160.94，按 2026Q1 归母净利和 EPS 反推的市值约 CNY1606 亿元。2026Q1 收入 CNY42.66 亿元，归母净利 CNY6.38 亿元，毛利率 19.8\%，经营现金流 CNY-11.27 亿元。以上是估值分母，不允许被行业叙事替代。

估值上，本报告采用 \textbf{PE/PB/PS} 方法：70\% PE + 10\% PB + 20\% PS，用于有正 EPS 但 AI 纯度不完美的模块/器件平台。华工科技 的 2026E/2027E/2028E EPS 分别为 CNY2.56、CNY3.02、CNY3.44，2027E 收入/股为 CNY20.18，BVPS 为 CNY11.62；内在组件权重为 PE 70\% / PB 10\% / PS 20\%，内在价值锚 CNY92.79，市场情绪锚 CNY128.75，综合权重 内在47\% / 市场38\% / 券商15\%，综合目标价 CNY128.75，综合空间 -20.0\%，动作 中性观察，证据质量 B。二级校验为：数通收入占比、毛利率桥、收入规模和现金转化。

下一轮调研首先验证催化剂：高速光模块收入占比提升、工业激光盈利改善。 其次验证失效条件：营运资本压力、混合业务折价扩大，或数通订单不放量。 若催化剂没有进入收入、毛利率和现金流，不能上调 PE；若失效条件出现，应同时下调 EPS 和动作标签。对于产业链映射，必须核查它的增长是否来自本层级真实需求，而不是来自板块情绪扩散。


\subsection{剑桥科技（603083）}
剑桥科技 的产业定位是 \textbf{光模块与通信设备高 beta 标的}，在本报告中的层级为 高弹性主题，研究权重 2\%。当前价 CNY245.92，按 2026Q1 归母净利和 EPS 反推的市值约 CNY856 亿元。2026Q1 收入 CNY12.87 亿元，归母净利 CNY1.18 亿元，毛利率 29.3\%，经营现金流 CNY-2.19 亿元。以上是估值分母，不允许被行业叙事替代。

估值上，本报告采用 \textbf{PE/PB/PS} 方法：45\% PE + 15\% PB + 40\% PS，用于小规模但具战略稀缺性的光芯片、相干模块和精密光学。剑桥科技 的 2026E/2027E/2028E EPS 分别为 CNY1.70、CNY2.21、CNY2.76，2027E 收入/股为 CNY24.03，BVPS 为 CNY21.45；内在组件权重为 PE 45\% / PB 15\% / PS 40\%，内在价值锚 CNY105.68，市场情绪锚 CNY177.06，综合权重 内在47\% / 市场38\% / 券商15\%，综合目标价 CNY177.06，综合空间 -28.0\%，动作 中性观察，证据质量 C+。二级校验为：客户认证、收入规模、毛利率持续性和战略稀缺性。

下一轮调研首先验证催化剂：800G/1.6T 订单持续转化、费用率被收入放量摊薄。 其次验证失效条件：Q2/Q3 利润不放量、现金流仍为负，或客户结构恶化。 若催化剂没有进入收入、毛利率和现金流，不能上调 PE；若失效条件出现，应同时下调 EPS 和动作标签。对于产业链映射，必须核查它的增长是否来自本层级真实需求，而不是来自板块情绪扩散。


\subsection{太辰光（300570）}
太辰光 的产业定位是 \textbf{光纤连接器与无源光器件}，在本报告中的层级为 无源器件，研究权重 2\%。当前价 CNY234.10，按 2026Q1 归母净利和 EPS 反推的市值约 CNY531 亿元。2026Q1 收入 CNY3.41 亿元，归母净利 CNY0.66 亿元，毛利率 37.8\%，经营现金流 CNY-0.93 亿元。以上是估值分母，不允许被行业叙事替代。

估值上，本报告采用 \textbf{PE/PB/PS} 方法：45\% 周期正常化 PE + 30\% PB/ROE + 25\% PS，用于资产较重的光纤光缆和海缆平台。太辰光 的 2026E/2027E/2028E EPS 分别为 CNY1.16、CNY1.33、CNY1.49，2027E 收入/股为 CNY6.91，BVPS 为 CNY7.79；内在组件权重为 PE 45\% / PB 30\% / PS 25\%，内在价值锚 CNY26.60，市场情绪锚 CNY163.87，综合权重 内在38\% / 市场47\% / 券商15\%，综合目标价 CNY163.87，综合空间 -30.0\%，动作 中性观察，证据质量 C。二级校验为：周期正常化 EPS、经营现金流、运营商/项目订单、光纤光缆价格。

下一轮调研首先验证催化剂：AI 数据中心连接器订单恢复、Q2 利润重新加速。 其次验证失效条件：收入继续下滑、无源器件降价，或库存重建停滞。 若催化剂没有进入收入、毛利率和现金流，不能上调 PE；若失效条件出现，应同时下调 EPS 和动作标签。对于产业链映射，必须核查它的增长是否来自本层级真实需求，而不是来自板块情绪扩散。


\subsection{源杰科技（688498）}
源杰科技 的产业定位是 \textbf{激光芯片与高速光源}，在本报告中的层级为 战略上游，研究权重 4\%。当前价 CNY1776.91，按 2026Q1 归母净利和 EPS 反推的市值约 CNY1526 亿元。2026Q1 收入 CNY3.55 亿元，归母净利 CNY1.79 亿元，毛利率 77.8\%，经营现金流 CNY0.40 亿元。以上是估值分母，不允许被行业叙事替代。

估值上，本报告采用 \textbf{PE/PB/PS} 方法：45\% PE + 15\% PB + 40\% PS，用于小规模但具战略稀缺性的光芯片、相干模块和精密光学。源杰科技 的 2026E/2027E/2028E EPS 分别为 CNY9.09、CNY12.54、CNY16.30，2027E 收入/股为 CNY24.83，BVPS 为 CNY29.38；内在组件权重为 PE 45\% / PB 15\% / PS 40\%，内在价值锚 CNY434.08，市场情绪锚 CNY1421.53，综合权重 内在47\% / 市场38\% / 券商15\%，综合目标价 CNY1421.53，综合空间 -20.0\%，动作 中性观察，证据质量 B。二级校验为：客户认证、收入规模、毛利率持续性和战略稀缺性。

下一轮调研首先验证催化剂：100G/200G EML 和硅光光源进入头部模块客户。 其次验证失效条件：客户认证延迟、良率不稳，或双供应削弱议价能力。 若催化剂没有进入收入、毛利率和现金流，不能上调 PE；若失效条件出现，应同时下调 EPS 和动作标签。对于产业链映射，必须核查它的增长是否来自本层级真实需求，而不是来自板块情绪扩散。


\subsection{亨通光电（600487）}
亨通光电 的产业定位是 \textbf{光纤光缆、海缆与电力通信一体化}，在本报告中的层级为 光纤光缆，研究权重 4\%。当前价 CNY110.81，按 2026Q1 归母净利和 EPS 反推的市值约 CNY2684 亿元。2026Q1 收入 CNY177.91 亿元，归母净利 CNY11.05 亿元，毛利率 16.0\%，经营现金流 CNY-1.94 亿元。以上是估值分母，不允许被行业叙事替代。

估值上，本报告采用 \textbf{PE/PB/PS} 方法：45\% 周期正常化 PE + 30\% PB/ROE + 25\% PS，用于资产较重的光纤光缆和海缆平台。亨通光电 的 2026E/2027E/2028E EPS 分别为 CNY1.90、CNY2.13、CNY2.34，2027E 收入/股为 CNY34.27，BVPS 为 CNY13.14；内在组件权重为 PE 45\% / PB 30\% / PS 25\%，内在价值锚 CNY36.95，市场情绪锚 CNY86.43，综合权重 内在38\% / 市场47\% / 券商15\%，综合目标价 CNY86.43，综合空间 -22.0\%，动作 中性观察，证据质量 B。二级校验为：周期正常化 EPS、经营现金流、运营商/项目订单、光纤光缆价格。

下一轮调研首先验证催化剂：运营商/云侧光纤需求修复、海缆订单、营运资本改善。 其次验证失效条件：电信资本开支走弱、光缆价格承压，或现金转化落后收入。 若催化剂没有进入收入、毛利率和现金流，不能上调 PE；若失效条件出现，应同时下调 EPS 和动作标签。对于产业链映射，必须核查它的增长是否来自本层级真实需求，而不是来自板块情绪扩散。


\subsection{中天科技（600522）}
中天科技 的产业定位是 \textbf{光纤光缆、海缆与通信线缆}，在本报告中的层级为 光纤光缆，研究权重 4\%。当前价 CNY61.29，按 2026Q1 归母净利和 EPS 反推的市值约 CNY2071 亿元。2026Q1 收入 CNY131.42 亿元，归母净利 CNY9.19 亿元，毛利率 15.6\%，经营现金流 CNY-19.49 亿元。以上是估值分母，不允许被行业叙事替代。

估值上，本报告采用 \textbf{PE/PB/PS} 方法：45\% 周期正常化 PE + 30\% PB/ROE + 25\% PS，用于资产较重的光纤光缆和海缆平台。中天科技 的 2026E/2027E/2028E EPS 分别为 CNY1.13、CNY1.27、CNY1.40，2027E 收入/股为 CNY18.15，BVPS 为 CNY11.14；内在组件权重为 PE 45\% / PB 30\% / PS 25\%，内在价值锚 CNY23.35，市场情绪锚 CNY47.81，综合权重 内在38\% / 市场47\% / 券商15\%，综合目标价 CNY47.81，综合空间 -22.0\%，动作 中性观察，证据质量 B。二级校验为：周期正常化 EPS、经营现金流、运营商/项目订单、光纤光缆价格。

下一轮调研首先验证催化剂：海缆交付、运营商光纤需求改善、毛利率修复。 其次验证失效条件：电网/海缆项目延迟、光纤价格下降，或经营现金流恶化。 若催化剂没有进入收入、毛利率和现金流，不能上调 PE；若失效条件出现，应同时下调 EPS 和动作标签。对于产业链映射，必须核查它的增长是否来自本层级真实需求，而不是来自板块情绪扩散。


\subsection{长飞光纤（601869）}
长飞光纤 的产业定位是 \textbf{光纤预制棒、光纤与光缆}，在本报告中的层级为 光纤/预制棒，研究权重 4\%。当前价 CNY543.00，按 2026Q1 归母净利和 EPS 反推的市值约 CNY4481 亿元。2026Q1 收入 CNY36.95 亿元，归母净利 CNY4.95 亿元，毛利率 41.5\%，经营现金流 CNY6.18 亿元。以上是估值分母，不允许被行业叙事替代。

估值上，本报告采用 \textbf{PE/PB/PS} 方法：45\% 周期正常化 PE + 30\% PB/ROE + 25\% PS，用于资产较重的光纤光缆和海缆平台。长飞光纤 的 2026E/2027E/2028E EPS 分别为 CNY2.61、CNY3.03、CNY3.39，2027E 收入/股为 CNY22.58，BVPS 为 CNY17.19；内在组件权重为 PE 45\% / PB 30\% / PS 25\%，内在价值锚 CNY51.93，市场情绪锚 CNY434.40，综合权重 内在38\% / 市场47\% / 券商15\%，综合目标价 CNY434.40，综合空间 -20.0\%，动作 中性观察，证据质量 B+。二级校验为：周期正常化 EPS、经营现金流、运营商/项目订单、光纤光缆价格。

下一轮调研首先验证催化剂：预制棒利用率修复、云/园区光纤建设、海外订单。 其次验证失效条件：光纤供给过剩、ASP 下滑，或 Q2/Q3 利润未延续高 Q1 利润率。 若催化剂没有进入收入、毛利率和现金流，不能上调 PE；若失效条件出现，应同时下调 EPS 和动作标签。对于产业链映射，必须核查它的增长是否来自本层级真实需求，而不是来自板块情绪扩散。


\subsection{烽火通信（600498）}
烽火通信 的产业定位是 \textbf{电信光网络设备与系统}，在本报告中的层级为 网络设备，研究权重 3\%。当前价 CNY76.22，按 2026Q1 归母净利和 EPS 反推的市值约 CNY975 亿元。2026Q1 收入 CNY45.10 亿元，归母净利 CNY0.38 亿元，毛利率 26.4\%，经营现金流 CNY-7.88 亿元。以上是估值分母，不允许被行业叙事替代。

估值上，本报告采用 \textbf{PE/PB/PS} 方法：55\% PE + 20\% PB + 25\% PS，用于运营商/云网络设备业务。烽火通信 的 2026E/2027E/2028E EPS 分别为 CNY0.14、CNY0.16、CNY0.18，2027E 收入/股为 CNY18.42，BVPS 为 CNY12.94；内在组件权重为 PE 55\% / PB 20\% / PS 25\%，内在价值锚 CNY12.45，市场情绪锚 CNY53.35，综合权重 内在45\% / 市场55\% / 券商0\%，综合目标价 CNY53.35，综合空间 -30.0\%，动作 中性观察，证据质量 B-。二级校验为：在手订单、运营商资本开支、业务结构、现金流和非光通信业务稀释。

下一轮调研首先验证催化剂：OTN/ROADM 需求、运营商骨干网升级、盈利正常化。 其次验证失效条件：运营商资本开支延迟、低毛利项目占比提高，或净利率未修复。 若催化剂没有进入收入、毛利率和现金流，不能上调 PE；若失效条件出现，应同时下调 EPS 和动作标签。对于产业链映射，必须核查它的增长是否来自本层级真实需求，而不是来自板块情绪扩散。


\subsection{中兴通讯（000063）}
中兴通讯 的产业定位是 \textbf{网络设备、运营商/云基础设施与终端应用}，在本报告中的层级为 网络设备/应用，研究权重 6\%。当前价 CNY35.67，按 2026Q1 归母净利和 EPS 反推的市值约 CNY1731 亿元。2026Q1 收入 CNY349.88 亿元，归母净利 CNY13.10 亿元，毛利率 28.3\%，经营现金流 CNY-19.79 亿元。以上是估值分母，不允许被行业叙事替代。

估值上，本报告采用 \textbf{PE/PB/PS} 方法：55\% PE + 20\% PB + 25\% PS，用于运营商/云网络设备业务。中兴通讯 的 2026E/2027E/2028E EPS 分别为 CNY1.23、CNY1.37、CNY1.51，2027E 收入/股为 CNY36.70，BVPS 为 CNY16.03；内在组件权重为 PE 55\% / PB 20\% / PS 25\%，内在价值锚 CNY31.03，市场情绪锚 CNY25.68，综合权重 内在50\% / 市场35\% / 券商15\%，综合目标价 CNY31.06，综合空间 -12.9\%，动作 中性，证据质量 A-。二级校验为：在手订单、运营商资本开支、业务结构、现金流和非光通信业务稀释。

下一轮调研首先验证催化剂：运营商网络升级、AI 服务器/网络产品进展、利润率改善。 其次验证失效条件：运营商资本开支走弱、海外政策压力，或非光通信业务稀释盈利质量。 若催化剂没有进入收入、毛利率和现金流，不能上调 PE；若失效条件出现，应同时下调 EPS 和动作标签。对于产业链映射，必须核查它的增长是否来自本层级真实需求，而不是来自板块情绪扩散。


\subsection{仕佳光子（688313）}
仕佳光子 的产业定位是 \textbf{PLC 分路器、AWG 与无源光芯片/器件}，在本报告中的层级为 无源光芯片，研究权重 3\%。当前价 CNY179.70，按 2026Q1 归母净利和 EPS 反推的市值约 CNY812 亿元。2026Q1 收入 CNY5.77 亿元，归母净利 CNY1.16 亿元，毛利率 34.1\%，经营现金流 CNY-0.97 亿元。以上是估值分母，不允许被行业叙事替代。

估值上，本报告采用 \textbf{PE/PB/PS} 方法：45\% PE + 15\% PB + 40\% PS，用于小规模但具战略稀缺性的光芯片、相干模块和精密光学。仕佳光子 的 2026E/2027E/2028E EPS 分别为 CNY1.12、CNY1.40、CNY1.65，2027E 收入/股为 CNY6.93，BVPS 为 CNY3.67；内在组件权重为 PE 45\% / PB 15\% / PS 40\%，内在价值锚 CNY47.49，市场情绪锚 CNY125.79，综合权重 内在47\% / 市场38\% / 券商15\%，综合目标价 CNY125.79，综合空间 -30.0\%，动作 中性观察，证据质量 B。二级校验为：客户认证、收入规模、毛利率持续性和战略稀缺性。

下一轮调研首先验证催化剂：AWG/PLC 器件在高速模块和电信/数据中心 WDM 系统中放量。 其次验证失效条件：光芯片降价、客户认证延迟，或产能利用率低于 Q1 水平。 若催化剂没有进入收入、毛利率和现金流，不能上调 PE；若失效条件出现，应同时下调 EPS 和动作标签。对于产业链映射，必须核查它的增长是否来自本层级真实需求，而不是来自板块情绪扩散。


\subsection{长光华芯（688048）}
长光华芯 的产业定位是 \textbf{高功率与数通激光芯片}，在本报告中的层级为 激光芯片，研究权重 2\%。当前价 CNY396.00，按 2026Q1 归母净利和 EPS 反推的市值约 CNY698 亿元。2026Q1 收入 CNY1.30 亿元，归母净利 CNY0.04 亿元，毛利率 30.1\%，经营现金流 CNY-1.31 亿元。以上是估值分母，不允许被行业叙事替代。

估值上，本报告采用 \textbf{PE/PB/PS} 方法：45\% PE + 15\% PB + 40\% PS，用于小规模但具战略稀缺性的光芯片、相干模块和精密光学。长光华芯 的 2026E/2027E/2028E EPS 分别为 CNY0.13、CNY0.18、CNY0.25，2027E 收入/股为 CNY5.34，BVPS 为 CNY17.10；内在组件权重为 PE 45\% / PB 15\% / PS 40\%，内在价值锚 CNY26.74，市场情绪锚 CNY277.20，综合权重 内在47\% / 市场38\% / 券商15\%，综合目标价 CNY277.20，综合空间 -30.0\%，动作 中性观察，证据质量 C+。二级校验为：客户认证、收入规模、毛利率持续性和战略稀缺性。

下一轮调研首先验证催化剂：数通激光芯片认证、良率改善、高功率激光盈利修复。 其次验证失效条件：研发转化延迟、良率不稳，或 Q2/Q3 盈利不能验证高倍数。 若催化剂没有进入收入、毛利率和现金流，不能上调 PE；若失效条件出现，应同时下调 EPS 和动作标签。对于产业链映射，必须核查它的增长是否来自本层级真实需求，而不是来自板块情绪扩散。


\subsection{光库科技（300620）}
光库科技 的产业定位是 \textbf{光纤器件、隔离器与铌酸锂/调制器期权}，在本报告中的层级为 光器件，研究权重 3\%。当前价 CNY362.43，按 2026Q1 归母净利和 EPS 反推的市值约 CNY903 亿元。2026Q1 收入 CNY4.26 亿元，归母净利 CNY0.45 亿元，毛利率 36.6\%，经营现金流 CNY-0.34 亿元。以上是估值分母，不允许被行业叙事替代。

估值上，本报告采用 \textbf{PE/PB/PS} 方法：45\% 周期正常化 PE + 30\% PB/ROE + 25\% PS，用于资产较重的光纤光缆和海缆平台。光库科技 的 2026E/2027E/2028E EPS 分别为 CNY0.78、CNY0.98、CNY1.15，2027E 收入/股为 CNY9.30，BVPS 为 CNY8.85；内在组件权重为 PE 45\% / PB 30\% / PS 25\%，内在价值锚 CNY28.59，市场情绪锚 CNY217.46，综合权重 内在50\% / 市场35\% / 券商15\%，综合目标价 CNY116.95，综合空间 -67.7\%，动作 减持，证据质量 B。二级校验为：周期正常化 EPS、经营现金流、运营商/项目订单、光纤光缆价格。

下一轮调研首先验证催化剂：高速光器件附加值提升、薄膜铌酸锂商业化、海外需求修复。 其次验证失效条件：调制器商业化延迟、毛利率压缩，或 AI 相关器件需求低于预期。 若催化剂没有进入收入、毛利率和现金流，不能上调 PE；若失效条件出现，应同时下调 EPS 和动作标签。对于产业链映射，必须核查它的增长是否来自本层级真实需求，而不是来自板块情绪扩散。


\subsection{通鼎互联（002491）}
通鼎互联 的产业定位是 \textbf{光纤光缆、通信线缆与网络集成}，在本报告中的层级为 光纤光缆，研究权重 3\%。当前价 CNY36.15，按 2026Q1 归母净利和 EPS 反推的市值约 CNY444 亿元。2026Q1 收入 CNY10.64 亿元，归母净利 CNY0.51 亿元，毛利率 27.0\%，经营现金流 CNY-0.18 亿元。以上是估值分母，不允许被行业叙事替代。

估值上，本报告采用 \textbf{PE/PB/PS} 方法：20\% 周期正常化 PE + 30\% PB/ROE + 50\% 收入期权 PS，用于光纤光缆和通信集成敞口。通鼎互联 的 2026E/2027E/2028E EPS 分别为 CNY0.18、CNY0.20、CNY0.22，2027E 收入/股为 CNY4.21，BVPS 为 CNY2.05；内在组件权重为 PE 20\% / PB 30\% / PS 50\%，内在价值锚 CNY13.96，市场情绪锚 CNY26.03，综合权重 内在45\% / 市场55\% / 券商0\%，综合目标价 CNY26.03，综合空间 -28.0\%，动作 中性观察，证据质量 B-。二级校验为：现金转化、运营商/云侧线缆订单、线缆 ASP、AI/数据中心敞口纯度。

下一轮调研首先验证催化剂：电信光纤/线缆需求恢复、通信线缆订单改善、现金流修复。 其次验证失效条件：运营商资本开支延迟、线缆 ASP 承压，或 Q2/Q3 利润不能确认 Q1 修复。 若催化剂没有进入收入、毛利率和现金流，不能上调 PE；若失效条件出现，应同时下调 EPS 和动作标签。对于产业链映射，必须核查它的增长是否来自本层级真实需求，而不是来自板块情绪扩散。


\subsection{永鼎股份（600105）}
永鼎股份 的产业定位是 \textbf{光缆、通信线缆与电信基础设施集成}，在本报告中的层级为 光纤光缆，研究权重 3\%。当前价 CNY65.71，按 2026Q1 归母净利和 EPS 反推的市值约 CNY961 亿元。2026Q1 收入 CNY12.46 亿元，归母净利 CNY1.59 亿元，毛利率 26.2\%，经营现金流 CNY2.59 亿元。以上是估值分母，不允许被行业叙事替代。

估值上，本报告采用 \textbf{PE/PB/PS} 方法：20\% 周期正常化 PE + 30\% PB/ROE + 50\% 收入期权 PS，用于光纤光缆和通信集成敞口。永鼎股份 的 2026E/2027E/2028E EPS 分别为 CNY0.45、CNY0.50、CNY0.54，2027E 收入/股为 CNY3.90，BVPS 为 CNY2.27；内在组件权重为 PE 20\% / PB 30\% / PS 50\%，内在价值锚 CNY14.77，市场情绪锚 CNY46.00，综合权重 内在45\% / 市场55\% / 券商0\%，综合目标价 CNY46.00，综合空间 -30.0\%，动作 中性观察，证据质量 B。二级校验为：现金转化、运营商/云侧线缆订单、线缆 ASP、AI/数据中心敞口纯度。

下一轮调研首先验证催化剂：运营商线缆订单、电信基础设施交付、毛利率稳定。 其次验证失效条件：项目延迟、线缆价格下行，或营运资本压力抵消利润修复。 若催化剂没有进入收入、毛利率和现金流，不能上调 PE；若失效条件出现，应同时下调 EPS 和动作标签。对于产业链映射，必须核查它的增长是否来自本层级真实需求，而不是来自板块情绪扩散。


\subsection{长芯博创（300548）}
长芯博创 的产业定位是 \textbf{光模块、有源光器件与相干传输组件}，在本报告中的层级为 光模块/光器件，研究权重 4\%。当前价 CNY262.65，按 2026Q1 归母净利和 EPS 反推的市值约 CNY759 亿元。2026Q1 收入 CNY6.71 亿元，归母净利 CNY1.30 亿元，毛利率 49.2\%，经营现金流 CNY2.26 亿元。以上是估值分母，不允许被行业叙事替代。

估值上，本报告采用 \textbf{PE/PB/PS} 方法：70\% PE + 10\% PB + 20\% PS，用于有正 EPS 但 AI 纯度不完美的模块/器件平台。长芯博创 的 2026E/2027E/2028E EPS 分别为 CNY1.88、CNY2.34、CNY2.77，2027E 收入/股为 CNY12.08，BVPS 为 CNY7.17；内在组件权重为 PE 70\% / PB 10\% / PS 20\%，内在价值锚 CNY80.23，市场情绪锚 CNY189.11，综合权重 内在55\% / 市场45\% / 券商0\%，综合目标价 CNY189.11，综合空间 -28.0\%，动作 中性观察，证据质量 B。二级校验为：数通收入占比、毛利率桥、收入规模和现金转化。

下一轮调研首先验证催化剂：数通模块订单、相干产品增长、毛利率维持。 其次验证失效条件：高速产品爬坡延迟、模块 ASP 承压，或毛利率回落。 若催化剂没有进入收入、毛利率和现金流，不能上调 PE；若失效条件出现，应同时下调 EPS 和动作标签。对于产业链映射，必须核查它的增长是否来自本层级真实需求，而不是来自板块情绪扩散。


\subsection{德科立（688205）}
德科立 的产业定位是 \textbf{相干光模块与传输设备组件}，在本报告中的层级为 相干/传输，研究权重 3\%。当前价 CNY210.52，按 2026Q1 归母净利和 EPS 反推的市值约 CNY345 亿元。2026Q1 收入 CNY2.54 亿元，归母净利 CNY0.20 亿元，毛利率 25.7\%，经营现金流 CNY-0.64 亿元。以上是估值分母，不允许被行业叙事替代。

估值上，本报告采用 \textbf{PE/PB/PS} 方法：45\% PE + 15\% PB + 40\% PS，用于小规模但具战略稀缺性的光芯片、相干模块和精密光学。德科立 的 2026E/2027E/2028E EPS 分别为 CNY0.55、CNY0.68、CNY0.80，2027E 收入/股为 CNY8.76，BVPS 为 CNY14.79；内在组件权重为 PE 45\% / PB 15\% / PS 40\%，内在价值锚 CNY40.10，市场情绪锚 CNY134.73，综合权重 内在65\% / 市场35\% / 券商0\%，综合目标价 CNY73.50，综合空间 -65.1\%，动作 减持，证据质量 B。二级校验为：客户认证、收入规模、毛利率持续性和战略稀缺性。

下一轮调研首先验证催化剂：DCI/相干需求、海外客户拓展、盈利正常化。 其次验证失效条件：相干订单延迟、客户集中度压力，或净利率不能放大。 若催化剂没有进入收入、毛利率和现金流，不能上调 PE；若失效条件出现，应同时下调 EPS 和动作标签。对于产业链映射，必须核查它的增长是否来自本层级真实需求，而不是来自板块情绪扩散。


\subsection{腾景科技（688195）}
腾景科技 的产业定位是 \textbf{光通信和激光用精密光学元件}，在本报告中的层级为 精密光学，研究权重 3\%。当前价 CNY187.60，按 2026Q1 归母净利和 EPS 反推的市值约 CNY246 亿元。2026Q1 收入 CNY1.71 亿元，归母净利 CNY0.14 亿元，毛利率 37.2\%，经营现金流 CNY-0.05 亿元。以上是估值分母，不允许被行业叙事替代。

估值上，本报告采用 \textbf{PE/PB/PS} 方法：45\% PE + 15\% PB + 40\% PS，用于小规模但具战略稀缺性的光芯片、相干模块和精密光学。腾景科技 的 2026E/2027E/2028E EPS 分别为 CNY0.48、CNY0.57、CNY0.66，2027E 收入/股为 CNY6.81，BVPS 为 CNY7.87；内在组件权重为 PE 45\% / PB 15\% / PS 40\%，内在价值锚 CNY29.18，市场情绪锚 CNY112.56，综合权重 内在55\% / 市场30\% / 券商15\%，综合目标价 CNY83.18，综合空间 -55.7\%，动作 减持，证据质量 B。二级校验为：客户认证、收入规模、毛利率持续性和战略稀缺性。

下一轮调研首先验证催化剂：高端光学元件订单、激光/通信需求、毛利率改善。 其次验证失效条件：订单波动、收入规模不足，或光学元件价格承压。 若催化剂没有进入收入、毛利率和现金流，不能上调 PE；若失效条件出现，应同时下调 EPS 和动作标签。对于产业链映射，必须核查它的增长是否来自本层级真实需求，而不是来自板块情绪扩散。


\subsection{联特科技（301205）}
联特科技 的产业定位是 \textbf{高速光模块与数通收发器}，在本报告中的层级为 高弹性光模块，研究权重 2\%。当前价 CNY315.86，按 2026Q1 归母净利和 EPS 反推的市值约 CNY409 亿元。2026Q1 收入 CNY2.13 亿元，归母净利 CNY0.03 亿元，毛利率 43.4\%，经营现金流 CNY-1.08 亿元。以上是估值分母，不允许被行业叙事替代。

估值上，本报告采用 \textbf{PE/PB/PS} 方法：45\% PE + 15\% PB + 40\% PS，用于小规模但具战略稀缺性的光芯片、相干模块和精密光学。联特科技 的 2026E/2027E/2028E EPS 分别为 CNY0.12、CNY0.16、CNY0.20，2027E 收入/股为 CNY11.11，BVPS 为 CNY12.55；内在组件权重为 PE 45\% / PB 15\% / PS 40\%，内在价值锚 CNY33.68，市场情绪锚 CNY189.52，综合权重 内在55\% / 市场30\% / 券商15\%，综合目标价 CNY116.17，综合空间 -63.2\%，动作 减持，证据质量 C+。二级校验为：客户认证、收入规模、毛利率持续性和战略稀缺性。

下一轮调研首先验证催化剂：高速数通订单转化、客户认证、毛利率扩张。 其次验证失效条件：认证延迟、Q2/Q3 利润不放量，或现金流不跟随收入。 若催化剂没有进入收入、毛利率和现金流，不能上调 PE；若失效条件出现，应同时下调 EPS 和动作标签。对于产业链映射，必须核查它的增长是否来自本层级真实需求，而不是来自板块情绪扩散。


\subsection{兆龙互连（300913）}
兆龙互连 的产业定位是 \textbf{高速数据线缆、通信线缆与铜互连}，在本报告中的层级为 高速线缆/互连，研究权重 2\%。当前价 CNY50.75，按 2026Q1 归母净利和 EPS 反推的市值约 CNY174 亿元。2026Q1 收入 CNY5.08 亿元，归母净利 CNY0.49 亿元，毛利率 21.1\%，经营现金流 CNY0.52 亿元。以上是估值分母，不允许被行业叙事替代。

估值上，本报告采用 \textbf{PE/PB/PS} 方法：40\% PE + 25\% PB + 35\% PS，用于连接器、线缆和高速互连混合业务。兆龙互连 的 2026E/2027E/2028E EPS 分别为 CNY0.59、CNY0.70、CNY0.79，2027E 收入/股为 CNY7.29，BVPS 为 CNY7.73；内在组件权重为 PE 40\% / PB 25\% / PS 35\%，内在价值锚 CNY21.10，市场情绪锚 CNY24.36，综合权重 内在59\% / 市场41\% / 券商0\%，综合目标价 CNY22.44，综合空间 -55.8\%，动作 减持，证据质量 B。二级校验为：AI 互连敞口纯度、产品结构、营运资本和毛利率稳定性。

下一轮调研首先验证催化剂：高速数据线缆需求、AI 服务器/网络互连订单、毛利率稳定。 其次验证失效条件：铜互连需求不及预期、产品结构恶化，或价格竞争加剧。 若催化剂没有进入收入、毛利率和现金流，不能上调 PE；若失效条件出现，应同时下调 EPS 和动作标签。对于产业链映射，必须核查它的增长是否来自本层级真实需求，而不是来自板块情绪扩散。


\subsection{意华股份（002897）}
意华股份 的产业定位是 \textbf{高速连接器与通信互连组件}，在本报告中的层级为 连接器/互连，研究权重 2\%。当前价 CNY84.50，按 2026Q1 归母净利和 EPS 反推的市值约 CNY162 亿元。2026Q1 收入 CNY12.92 亿元，归母净利 CNY0.40 亿元，毛利率 20.0\%，经营现金流 CNY2.74 亿元。以上是估值分母，不允许被行业叙事替代。

估值上，本报告采用 \textbf{PE/PB/PS} 方法：40\% PE + 25\% PB + 35\% PS，用于连接器、线缆和高速互连混合业务。意华股份 的 2026E/2027E/2028E EPS 分别为 CNY0.95、CNY1.10、CNY1.23，2027E 收入/股为 CNY35.23，BVPS 为 CNY14.66；内在组件权重为 PE 40\% / PB 25\% / PS 35\%，内在价值锚 CNY48.20，市场情绪锚 CNY52.39，综合权重 内在59\% / 市场41\% / 券商0\%，综合目标价 CNY49.92，综合空间 -40.9\%，动作 减持，证据质量 B-。二级校验为：AI 互连敞口纯度、产品结构、营运资本和毛利率稳定性。

下一轮调研首先验证催化剂：高速连接器订单、AI 服务器互连需求、盈利修复。 其次验证失效条件：连接器需求放缓、非光通信业务稀释，或毛利率压缩。 若催化剂没有进入收入、毛利率和现金流，不能上调 PE；若失效条件出现，应同时下调 EPS 和动作标签。对于产业链映射，必须核查它的增长是否来自本层级真实需求，而不是来自板块情绪扩散。


\subsection{神宇股份（300563）}
神宇股份 的产业定位是 \textbf{射频同轴线缆与高速通信互连}，在本报告中的层级为 线缆/互连，研究权重 2\%。当前价 CNY24.73，按 2026Q1 归母净利和 EPS 反推的市值约 CNY48 亿元。2026Q1 收入 CNY2.04 亿元，归母净利 CNY0.13 亿元，毛利率 18.2\%，经营现金流 CNY0.07 亿元。以上是估值分母，不允许被行业叙事替代。

估值上，本报告采用 \textbf{PE/PB/PS} 方法：40\% PE + 25\% PB + 35\% PS，用于连接器、线缆和高速互连混合业务。神宇股份 的 2026E/2027E/2028E EPS 分别为 CNY0.32、CNY0.37、CNY0.41，2027E 收入/股为 CNY5.52，BVPS 为 CNY6.33；内在组件权重为 PE 40\% / PB 25\% / PS 35\%，内在价值锚 CNY13.70，市场情绪锚 CNY11.87，综合权重 内在59\% / 市场41\% / 券商0\%，综合目标价 CNY12.95，综合空间 -47.7\%，动作 减持，证据质量 C+。二级校验为：AI 互连敞口纯度、产品结构、营运资本和毛利率稳定性。

下一轮调研首先验证催化剂：高速通信线缆需求、客户结构改善、毛利率稳定。 其次验证失效条件：低毛利产品占比提高、订单转化弱，或现金流落后利润。 若催化剂没有进入收入、毛利率和现金流，不能上调 PE；若失效条件出现，应同时下调 EPS 和动作标签。对于产业链映射，必须核查它的增长是否来自本层级真实需求，而不是来自板块情绪扩散。
